\section{Estrategia de lanzamiento: de círculos reducidos al efecto de red}

La difusión inicial de Elys provino de círculos reducidos, en particular de personas atentas a los nuevos paradigmas de la IA. Este modo de arranque tiene ventajas: los primeros usuarios entienden los conceptos de context, doble digital y proactive, y están dispuestos a tolerar un producto todavía tosco. La desventaja es que un círculo reducido no equivale al mercado masivo; al salir de ese círculo, los problemas de privacidad, los malentendidos, el abuso y las interacciones de baja calidad se amplifican.

\subsection{Ruta de lanzamiento en tres etapas}

\noindent\begin{tabularx}{\textwidth}{p{0.22\textwidth}p{0.34\textwidth}X}
\toprule
Etapa & Objetivo & Riesgos principales \\
\midrule
Círculo semilla & Validar el paradigma y el volante de inercia central & Los usuarios saben demasiado de IA y no representan al público general. \\
Comunidades verticales & Validar la tasa de conexión entre personas reales y la calidad de las relaciones & Aumenta la presión sobre los límites de la comunidad, la privacidad y la gobernanza. \\
Mercado masivo & Ampliar el efecto de red y la monetización & Ruido de IA, abuso, confianza y riesgos regulatorios. \\
\bottomrule
\end{tabularx}

\begin{importantbox}{Salir del círculo inicial no es solo un problema de crecimiento}
Cuando una red social de IA sale de su círculo inicial, cambia la estructura de riesgos. Los primeros usuarios tratan al doble digital como un experimento; los usuarios del público general pueden tomarlo por la persona misma. El producto debe establecer mecanismos de permisos, señalización, apelación y gobernanza antes de salir del círculo inicial.
\end{importantbox}

\subsection{Resumen de la sección}

Un producto como Elys no puede limitarse a buscar un crecimiento rápido. La estrategia de lanzamiento de una red social de IA debe construirse a la par de los mecanismos de confianza; de lo contrario, el efecto de red crecerá junto con los riesgos.
